\documentclass[aspectratio=169,10pt]{beamer}
\usetheme[numbering=fraction,progressbar=frametitle]{metropolis}
\usepackage{fontspec}
\setsansfont{Latin Modern Sans}
\usepackage{amsmath,amssymb,mathtools,graphicx,booktabs}
\usepackage{polyglossia}
\setdefaultlanguage{french}
\setbeamertemplate{navigation symbols}{}
\setbeamerfont{frametitle}{size=\large}
\setbeamerfont{normal text}{size=\normalsize}
\setlength{\parskip}{.45em}
\hypersetup{pdfauthor={},pdfsubject={Lecture slides},pdfkeywords={},pdftitle={Nombres complexes : point de vue algébrique}}
\providecommand{\N}{\mathbb{N}}
\providecommand{\Z}{\mathbb{Z}}
\providecommand{\D}{\mathbb{D}}
\providecommand{\Q}{\mathbb{Q}}
\providecommand{\R}{\mathbb{R}}
\providecommand{\C}{\mathbb{C}}
\providecommand{\vect}[1]{\overrightarrow{#1}}
\providecommand{\vi}{\vec{\imath}}
\providecommand{\vj}{\vec{\jmath}}
\providecommand{\vk}{\vec{k}}
\providecommand{\norm}[1]{\left\lVert #1\right\rVert}
\providecommand{\vc}[2]{\begin{pmatrix}#1\\#2\end{pmatrix}}
\providecommand{\deter}[2]{\det\left(#1,#2\right)}
\providecommand{\abs}[1]{\left\lvert #1\right\rvert}
\providecommand{\dd}{\,\mathrm{d}}
\providecommand{\e}{\mathrm{e}}
\providecommand{\Card}{\operatorname{Card}}
\providecommand{\E}{\operatorname{E}}
\providecommand{\V}{\operatorname{V}}
\providecommand{\degre}{^{\circ}}
\providecommand{\un}[1]{\,\mathrm{#1}}
\providecommand{\sci}[2]{#1\times10^{#2}}
\providecommand{\pgcd}{\operatorname{PGCD}}
\providecommand{\ppcm}{\operatorname{PPCM}}
\title{Nombres complexes : point de vue algébrique}
\subtitle{Cours et exemples}
\date{}
\author{}
\begin{document}
\begin{frame}[plain]
\vfill
{\usebeamerfont{title}\inserttitle\par}
\medskip
{\usebeamerfont{subtitle}\insertsubtitle\par}
\medskip
{\usebeamercolor[fg]{progress bar}\rule{\textwidth}{.5pt}}
\vfill
\end{frame}
\begin{frame}{Objectifs}
Calculer avec la forme algébrique ; utiliser le conjugué pour obtenir un inverse ou un quotient ; résoudre une équation du premier degré, ou une équation simple faisant intervenir \(z\) et \(\overline{z}\) ; développer une puissance par la formule du binôme.
\end{frame}
\begin{frame}{Activité d'introduction}
Au xvie siècle, on dispose d'une formule pour résoudre les équations du type \(x^3=px+q\). Elle est donnée ici et admise :
\[x=\sqrt[3]{\dfrac q2+\sqrt{\dfrac{q^2}{4}-\dfrac{p^3}{27}}}+\sqrt[3]{\dfrac q2-\sqrt{\dfrac{q^2}{4}-\dfrac{p^3}{27}}}.\]
On étudie l'équation \(x^3=15x+4\).
\end{frame}
\begin{frame}{Activité d'introduction (2)}
\begin{enumerate}
\setcounter{enumi}{0}\item Vérifier que \(4\) est solution de l'équation.

\item Calculer \(\dfrac{q^2}{4}-\dfrac{p^3}{27}\) pour \(p=15\) et \(q=4\). Que devient la formule ? Peut-on l'appliquer dans \(\R\) ?

\item Un calculateur de l'époque décide de continuer comme si un nombre noté \(\sqrt{-1}\) existait, avec la seule règle \(\sqrt{-1}\times\sqrt{-1}=-1\), et les règles de calcul habituelles pour le reste. Développer \(\left(2+\sqrt{-1}\right)^3\) avec ces règles.

\end{enumerate}
\end{frame}
\begin{frame}{Activité d'introduction (3)}
\begin{enumerate}
\setcounter{enumi}{3}\item Admettons que \(\sqrt{-121}\) s'écrive \(11\sqrt{-1}\). Montrer que le premier terme de la formule peut être pris égal à \(2+\sqrt{-1}\), puis trouver de même une valeur du second terme.

\item Que donne alors la formule ? Comparer avec la question 1.

\item Ce calcul est-il une démonstration ? Quelles règles faudrait-il fixer pour qu'il le devienne ?

\end{enumerate}
\end{frame}
\begin{frame}{L'ensemble des nombres complexes}
On admet l'existence d'un ensemble noté \(\C\), appelé ensemble des \emph{nombres complexes}, tel que :

\begin{enumerate}
\setcounter{enumi}{0}
\item \(\C\) contient \(\R\) ;

\item \(\C\) contient un nombre noté \(i\) tel que \(i^2=-1\) ;

\item \(\C\) est muni d'une addition et d'une multiplication qui prolongent celles de \(\R\) et suivent les mêmes règles de calcul (commutativité, associativité, distributivité) ;

\item tout nombre complexe \(z\) s'écrit de manière unique \(z=a+ib\) avec \(a\) et \(b\) réels.

\end{enumerate}
\end{frame}
\begin{frame}{L'ensemble des nombres complexes (2)}
L'écriture \(z=a+ib\) est la \emph{forme algébrique} de \(z\). Le réel \(a\) est la \emph{partie réelle} de \(z\), notée \(\operatorname{Re}(z)\) ; le réel \(b\) est la \emph{partie imaginaire} de \(z\), notée \(\operatorname{Im}(z)\). Si \(b=0\), \(z\) est réel ; si \(a=0\), \(z\) est \emph{imaginaire pur}.
\end{frame}
\begin{frame}{L'ensemble des nombres complexes}
Deux nombres complexes sont égaux si et seulement s'ils ont la même partie réelle et la même partie imaginaire. En particulier, \(a+ib=0\) avec \(a,b\) réels équivaut à \(a=b=0\).
\end{frame}
\begin{frame}{L'ensemble des nombres complexes}
La partie imaginaire de \(3-5i\) est \(-5\), et non \(-5i\) : c'est un réel. Il n'y a pas d'ordre compatible avec les opérations sur \(\C\) : on n'écrit jamais \(z<z'\) pour des complexes non réels.
\end{frame}
\begin{frame}{L'ensemble des nombres complexes}
\((2+3i)+(1-i)=3+2i\) et \((2+3i)(1-i)=2-2i+3i-3i^2=5+i\).

Les puissances de \(i\) se répètent de quatre en quatre : \(i^2=-1\), \(i^3=-i\), \(i^4=1\). Ainsi \(i^{4k+r}=i^r\) pour tout entier \(k\ge0\).
\end{frame}
\begin{frame}{Conjugué}
Le \emph{conjugué} de \(z=a+ib\) (\(a\), \(b\) réels) est \(\overline{z}=a-ib\).
\end{frame}
\begin{frame}{Conjugué}
Pour tout \(z=a+ib\) :
\[z+\overline z=2\operatorname{Re}(z),\qquad z-\overline z=2i\operatorname{Im}(z),\qquad z\overline z=a^2+b^2.\]
Ainsi \(z\) est réel si et seulement si \(\overline z=z\), et imaginaire pur si et seulement si \(\overline z=-z\). Le produit \(z\overline z\) est un réel positif, nul seulement pour \(z=0\).
\end{frame}
\begin{frame}{Conjugué}
\(z\overline z=(a+ib)(a-ib)=a^2-(ib)^2=a^2+b^2\). Les deux premières égalités s'obtiennent en ajoutant ou en soustrayant \(a+ib\) et \(a-ib\).
\end{frame}
\begin{frame}{Conjugué}
Tout nombre complexe \(z=a+ib\) non nul admet un inverse :
\[\dfrac1z=\dfrac{\overline z}{z\overline z}=\dfrac{a-ib}{a^2+b^2}.\]
\end{frame}
\begin{frame}{Conjugué}
On multiplie le numérateur et le dénominateur par le conjugué du dénominateur. Exemple : \(\dfrac{1}{3+4i}=\dfrac{3-4i}{9+16}=\dfrac{3}{25}-\dfrac{4}{25}i\).
\end{frame}
\begin{frame}{Conjugué}
Pour tous complexes \(z\), \(z'\) et tout entier naturel \(n\) :
\[\overline{z+z'}=\overline z+\overline{z'},\qquad\overline{zz'}=\overline z\,\overline{z'},\qquad\overline{z^n}=\overline z^{\,n},\]
et si \(z'\neq0\) : \(\overline{\left(\dfrac1{z'}\right)}=\dfrac1{\overline{z'}}\) et \(\overline{\left(\dfrac z{z'}\right)}=\dfrac{\overline z}{\overline{z'}}\).
\end{frame}
\begin{frame}{Conjugué}
\emph{Produit.} Posons \(z=a+ib\) et \(z'=a'+ib'\). Alors \(zz'=(aa'-bb')+i(ab'+a'b)\), donc \(\overline{zz'}=(aa'-bb')-i(ab'+a'b)\). D'autre part, \(\overline z\,\overline{z'}=(a-ib)(a'-ib')=(aa'-bb')-i(ab'+a'b)\). Les deux résultats sont égaux.
\end{frame}
\begin{frame}{Conjugué (2)}
\emph{Inverse.} Si \(z'\neq0\), \(z'\times\dfrac1{z'}=1\). Par la propriété du produit, \(\overline{z'}\times\overline{\left(\dfrac1{z'}\right)}=\overline1=1\), donc \(\overline{\left(\dfrac1{z'}\right)}\) est l'inverse de \(\overline{z'}\). Le quotient s'en déduit en écrivant \(\dfrac z{z'}=z\times\dfrac1{z'}\).
\end{frame}
\begin{frame}{Conjugué (3)}
\emph{Puissance.} Par récurrence sur \(n\). Pour \(n=0\) : \(\overline{z^0}=\overline1=1=\overline z^{\,0}\). Si \(\overline{z^n}=\overline z^{\,n}\) pour un entier \(n\), alors \(\overline{z^{n+1}}=\overline{z^n\times z}=\overline{z^n}\times\overline z=\overline z^{\,n}\times\overline z=\overline z^{\,n+1}\).
\end{frame}
\begin{frame}{Équations}
Si \(a\neq0\), l'équation \(az=b\) a une unique solution \(z=\dfrac ba\), que l'on écrit sous forme algébrique. Exemple : \((1+i)z=3-i\) donne \(z=\dfrac{(3-i)(1-i)}{2}=\dfrac{2-4i}{2}=1-2i\).
\end{frame}
\begin{frame}{Équations}
On pose \(z=x+iy\) avec \(x\), \(y\) réels, on écrit chaque membre sous forme algébrique, puis on identifie parties réelles et parties imaginaires. On obtient un système de deux équations réelles.
\end{frame}
\begin{frame}{Équations}
Avec \(z=x+iy\) : \(2x+2iy+i(x-iy)=(2x+y)+i(x+2y)\). L'identification donne \(2x+y=3\) et \(x+2y=-3\), d'où \(x=3\) et \(y=-3\). La solution est \(z=3-3i\). Vérification : \(2(3-3i)+i(3+3i)=6-6i+3i-3=3-3i\).
\end{frame}
\begin{frame}{Formule du binôme}
Pour tous complexes \(a\), \(b\) et tout entier naturel \(n\) :
\[(a+b)^n=\sum_{k=0}^{n}\binom nk a^kb^{n-k}.\]
\end{frame}
\begin{frame}{Formule du binôme}
Par récurrence sur \(n\). Pour \(n=0\), les deux membres valent \(1\). Supposons la formule vraie pour un entier \(n\). Alors
\[(a+b)^{n+1}=(a+b)\sum_{k=0}^{n}\binom nk a^kb^{n-k}=\sum_{k=0}^{n}\binom nk a^{k+1}b^{n-k}+\sum_{k=0}^{n}\binom nk a^kb^{n+1-k}.\]
Dans la première somme, on pose \(j=k+1\) : elle devient \(\sum_{j=1}^{n+1}\binom n{j-1}a^jb^{n+1-j}\). En regroupant les termes de même exposant de \(a\) et en utilisant la relation de Pascal \(\binom n{j-1}+\binom nj=\binom{n+1}j\) pour \(1\le j\le n\), on obtient
\[(a+b)^{n+1}=b^{n+1}+\sum_{j=1}^{n}\binom{n+1}j a^jb^{n+1-j}+a^{n+1}=\sum_{j=0}^{n+1}\binom{n+1}j a^jb^{n+1-j}.\]
La formule est vraie au rang \(n+1\), donc pour tout entier naturel \(n\).
\end{frame}
\begin{frame}{Formule du binôme}
\((1+i)^4=1+4i+6i^2+4i^3+i^4=1+4i-6-4i+1=-4\). On peut aussi écrire \((1+i)^2=2i\), puis \((1+i)^4=(2i)^2=-4\).
\end{frame}
\begin{frame}{Formule du binôme}
Forme algébrique unique ; \(i^2=-1\) ; \(z\overline z=a^2+b^2\) ; pour un quotient, multiplier par le conjugué du dénominateur ; pour une équation en \(z\) et \(\overline z\), poser \(z=x+iy\) et identifier.
\end{frame}
\end{document}
